
\chapter{Generics}

\section{Introduzione}
Prima dell'introduzione dei \textit{Generics}, la programmazione generica in Java veniva realizzata esclusivamente tramite l'ereditarietà. La classe \texttt{ArrayList}, ad esempio, gestiva semplicemente un array di riferimenti di tipo \texttt{Object}. Questo approccio presentava due problematiche critiche:

\begin{itemize}
    \item \textbf{Casting Obbligatorio}: Ogni volta che si recuperava un valore dalla lista, era necessario un cast esplicito al tipo desiderato.
    \begin{lstlisting}[language=Java]
ArrayList files = new ArrayList();
String filename = (String) files.get(0); // Cast necessario
    \end{lstlisting}
    \item \textbf{Assenza di Controllo sui Tipi}: Era possibile inserire oggetti di qualsiasi classe nella stessa collezione senza che il compilatore segnalasse errori.
    \begin{lstlisting}[language=Java]
files.add(new File("...")); // Compila correttamente, ma causera' un errore al runtime
    \end{lstlisting}
\end{itemize}
\medskip
I \textit{Generics} offrono una soluzione basata sui parametri di tipo. La sintassi permette di indicare esplicitamente il tipo di elementi che la collezione deve gestire.

\begin{itemize}
    \item \textbf{Sintassi Diamond}: Se il tipo e' dichiarato esplicitamente a sinistra, si può omettere nel costruttore usando l'operatore diamante \texttt{<>}.
    \begin{lstlisting}[language=Java]
ArrayList<String> files = new ArrayList<>();
    \end{lstlisting}
    \item \textbf{Inferenza con \texttt{var}}: Introdotta nelle versioni recenti, permette di dedurre il tipo dalla parte destra dell'assegnazione.
    \begin{lstlisting}[language=Java]
var files = new ArrayList<String>();
    \end{lstlisting}
\end{itemize}
\medskip
L'appello dei parametri di tipo risiede in due pilastri fondamentali:

\begin{enumerate}
    \item \textbf{Leggibilità}: È immediatamente chiaro quali oggetti siano contenuti nella lista (es. solo \texttt{String}).
    \item \textbf{Sicurezza (Type Safety)}: Il compilatore impedisce l'inserimento di oggetti di tipo errato.
\end{enumerate}

\begin{tip}[]
Un errore segnalato dal compilatore (\textit{compile-time error}) è di gran lunga preferibile a un'eccezione durante l'esecuzione (\textit{runtime ClassCastException}). I Generics spostano il controllo della correttezza del codice dalla fase di esecuzione alla fase di scrittura.
\end{tip}

La padronanza dei \textit{Generics} si sviluppa su tre livelli distinti, spostando l'attenzione dall'utilizzo alla progettazione.
\begin{enumerate}
    \item \textbf{Utilizzo (Livello Base)}: Uso delle classi fornite dal JDK (es. \texttt{ArrayList<String>}) come se fossero tipi nativi del linguaggio.
    \item \textbf{Risoluzione Problemi (Livello Intermedio)}: Comprensione dei meccanismi interni per gestire messaggi di errore complessi o interoperabilità con codice legacy.
    \item \textbf{Implementazione (Livello Avanzato)}: Creazione di proprie classi e metodi generici, anticipando i potenziali utilizzi futuri degli utenti.
\end{enumerate}
\medskip
Un compito cruciale del progettista è bilanciare sicurezza e usabilità. Consideriamo, ad esempio, il metodo \texttt{addAll}:
\begin{itemize}
    \item Deve essere possibile aggiungere \texttt{ArrayList<Manager>} a \texttt{ArrayList<Employee>} perche' un Manager \textit{è} un Employee.
    \item Non deve essere possibile il contrario.
\end{itemize}
\medskip
Vedremo presto che per risolvere queste asimmetrie, Java introduce i \textbf{wildcard types} (tipi jolly), che permettono di definire relazioni di sottotipo tra tipi parametrizzati.

\begin{tip}
La programmazione generica è particolarmente utile quando si scrive codice che altrimenti richiederebbe numerosi cast da tipi generali come \texttt{Object}. Per la maggior parte delle applicazioni comuni, le classi fornite dal JDK sono sufficienti.
\end{tip}

\section{Struttura di una Classe Generica}

Una classe generica è definita dall'introduzione di una o più \textbf{variabili di tipo} posizionate dopo il nome della classe. Queste fungono da segnaposto che verranno sostituiti dai tipi reali al momento dell'istanziazione.

\begin{lstlisting}[language=Java, caption={Esempio: la classe \texttt{Pair}}]
    public class Pair<T>{
        private T first;
        private T second;

        public Pair() { first = null; second = null; }
        public Pair(T first, T second) { this.first = first; 
                                         this.second = second; }

        public T getFirst() { return first; }
        public T getSecond() { return second; }

        public void setFirst(T newValue) { first = newValue; }
        public void setSecond(T newValue) { second = newValue; }
        }
\end{lstlisting}

La classe \texttt{Pair} dimostra come la variabile \texttt{T} sia utilizzata coerentemente in tutta la definizione:

\begin{itemize}
    \item \textbf{Campi}: \texttt{private T first;} --- Il tipo è determinato dall'utente.
    \item \textbf{Costruttore}: \texttt{public Pair(T first, T second)} --- Accetta parametri del tipo specificato.
    \item \textbf{Metodi}: \texttt{public T getFirst()} --- Restituisce un oggetto del tipo corretto senza necessità di cast.
\end{itemize}
\medskip
Per mantenere il codice leggibile e coerente con le API di Java, si utilizzano lettere singole maiuscole:
\begin{table}[ht]
\centering
\begin{tabular}{ll}
\hline
\textbf{Lettera} & \textbf{Significato} \\ \hline
\texttt{E} & Elemento (usato nelle collezioni) \\
\texttt{K} & Chiave (Key) \\
\texttt{V} & Valore (Value) \\
\texttt{T, U, S} & Tipi generici qualsiasi \\ \hline
\end{tabular}
\end{table}

È possibile definire classi con più variabili di tipo, ad esempio \texttt{public class Pair<T, U>}, permettendo alla classe di gestire coppie di oggetti di natura differente in modo \textit{type-safe}.

\section{Metodi Generici}

Un metodo generico può essere definito sia all'interno di classi ordinarie che di classi generiche. La sua caratteristica principale è la presenza di parametri di tipo indipendenti dalla classe.

La variabile di tipo deve essere dichiarata tra parentesi angolate \textit{dopo} i modificatori e \textit{prima} del tipo di ritorno:
\begin{center}
\texttt{modificatori <T> tipo\_ritorno nomeMetodo(parametri)}
\end{center}

Il compilatore Java è solitamente in grado di dedurre il tipo del parametro analizzando gli argomenti passati alla chiamata del metodo.
\begin{itemize}
    \item \textbf{Chiamata esplicita}: \texttt{ArrayAlg.<String>getMiddle("A", "B");}
    \item \textbf{Chiamata con inferenza}: \texttt{ArrayAlg.getMiddle("A", "B");}
\end{itemize}

\medskip

Quando vengono passati tipi eterogenei (es. \texttt{Double} e \texttt{Integer}), il compilatore cerca il \textit{common supertype}. Se esistono più supertipi validi (come \texttt{Number} e \texttt{Comparable}), l'inferenza fallisce.

\begin{tip}[\textbf{Il trucco dell'errore indotto}]
Per scoprire quale tipo il compilatore sta inferendo "sotto il cofano", prova ad assegnare il risultato del metodo a un tipo palesemente errato (es. \texttt{JButton}). L'errore del compilatore rivelerà l'esatta gerarchia di tipi (intersezione di tipi) che ha trovato, come \texttt{Object \& Serializable \& Comparable}.
\end{tip}

\section{Variabili di Tipo Vincolate (Bounded Type Variables)}

Talvolta è necessario limitare i tipi che possono essere sostituiti a una variabile di tipo \texttt{T}. Questo permette di invocare metodi specifici (come \texttt{compareTo}) sull'oggetto generico.

La clausola \texttt{<T extends BoundingType>} indica che \texttt{T} deve essere un \textbf{sottotipo} del tipo vincolante.
\begin{itemize}
    \item \textbf{Esempio}: \texttt{<T extends Comparable>} garantisce che \texttt{T} implementi l'interfaccia \texttt{Comparable}.
    \item \textbf{Nota Linguistica}: Si usa sempre \texttt{extends}, anche se il vincolo è un'interfaccia, per denotare il concetto generale di "sottotipo" senza aggiungere nuove keyword: in tale contesto, la parola chiave \texttt{implements} \textbf{non esiste}. 
\end{itemize}

\medskip

È possibile specificare più vincoli separandoli con l'operatore \texttt{\&} (la virgola è già usata per separare le variabili di tipo).
\begin{center}
\texttt{T extends Class1 \& Interface1 \& Interface2}
\end{center}

\begin{warningbox}[\textbf{Regole di Precedenza}]
\begin{enumerate}
    \item Si può specificare al massimo una classe come vincolo.
    \item Se presente, la classe deve essere indicata per \textbf{prima} nella lista dei vincoli.
    \item Le interfacce seguono la classe, separate da \texttt{\&}.
\end{enumerate}
\end{warningbox}

Ecco come appare il codice di un metodo \texttt{min} che calcola il minimo di un array di oggetti generici, purché siano comparabili:

\begin{lstlisting}[language=Java]
public static <T extends Comparable> T min(T[] a) {
    if (a == null || a.length == 0) return null;
    T smallest = a[0];
    for (int i = 1; i < a.length; i++) {
        if (smallest.compareTo(a[i]) > 0) {
            smallest = a[i];
        }
    }
    return smallest;
}
\end{lstlisting}

\begin{observation}
    L'interfaccia \texttt{Comparable} che abbiamo usato per l'esempio è essa stessa di tipo generico; ciò comporta che la sintassi che abbiamo usato genera dei warnings. Chiariremo questo dettaglio più avanti. 
\end{observation}

\section{Funzionamento interno dei Generics}
Vediamo ora come la JVM gestisce internamente i Generics. 

\subsection{Type Erasure (Cancellazione del Tipo)}

La \textit{Type Erasure} è il processo mediante il quale il compilatore Java traduce il codice generico in codice compatibile con le versioni precedenti della JVM (pre-Java 5).

Per ogni tipo generico, il compilatore genera un \textbf{Raw Type} (tipo grezzo) seguendo queste regole:
\begin{itemize}
    \item Le variabili di tipo senza vincoli (\texttt{<T>}) vengono sostituite da \texttt{Object}.
    \item Le variabili con vincoli (\texttt{<T extends B1 \& B2>}) vengono sostituite dal \textbf{primo vincolo} (\texttt{B1}).
\end{itemize}

\medskip

\begin{lstlisting}[language=Java, caption={Esempio: la classe \texttt{Interval}}]
    public class Interval<T extends Comparable & Serializable> implements Serializable{
        private T lower;
        private T upper;
        . . .
        public Interval(T first, T second)
        {
            if (first.compareTo(second) <= 0) { lower = first; upper = second; }
            else { lower = second; upper = first; }
        }
        }
\end{lstlisting}

Ad esempio, per la classe \texttt{Interval<T extends Comparable \& Serializable>}:
\begin{itemize}
    \item Il campo \texttt{private T lower} diventa \texttt{private Comparable lower}.
    \item Il metodo \texttt{compareTo} può essere chiamato direttamente senza cast.
\end{itemize}

\begin{warningbox}[\textbf{Efficienza e Ordine dei Vincoli}]
Se un tipo ha più vincoli, l'ordine non è indifferente. Il compilatore sostituisce \texttt{T} con il primo vincolo della lista. Se il primo vincolo è un'interfaccia "tag" (senza metodi, come \texttt{Serializable}), il compilatore dovrà inserire dei cast extra per invocare i metodi degli altri vincoli. 
Più in generale, per ridurre il \textit{runtime overhead}, conviene posizionare come primo vincolo l'interfaccia i cui metodi vengono invocati più frequentemente all'interno del codice generico.
\end{warningbox}

\subsection{Traduzione delle Espressioni Generiche}

Dopo che il compilatore ha eseguito la \textit{Type Erasure}, le chiamate ai metodi e l'accesso ai campi devono essere adattati per mantenere la coerenza dei tipi.
Poiché il tipo di ritorno di un metodo cancellato è solitamente \texttt{Object}, il compilatore inserisce automaticamente un \textbf{cast} nel codice chiamante. 
Data la chiamata \texttt{Employee buddy = buddies.getFirst();}, la traduzione nel bytecode avviene in due passaggi:

\begin{enumerate}
    \item Chiamata al metodo \textit{raw} \texttt{Pair.getFirst()} (ritorna \texttt{Object}).
    \item Conversione (\textit{cast}) del valore ritornato da \texttt{Object} a \texttt{Employee}.
\end{enumerate}

\medskip

Il medesimo principio si applica all'accesso diretto ai campi, qualora siano visibili. Anche se sintatticamente non appare alcun cast nel codice sorgente \texttt{.java}, il bytecode risultante conterrà un'istruzione di cast per garantire che il riferimento sia del tipo corretto.

Il vantaggio dei Generics non è quindi l'eliminazione dei cast a livello di esecuzione (che continuano a esistere nel bytecode), ma il fatto che il compilatore li gestisca \textbf{automaticamente} e in modo \textbf{sicuro}, sollevando il programmatore dal rischio di errori manuali. Praticamente, i Generics non sono altro che un modo di gestire in modo automatico ciò che prima della loro introduzione si faceva a mano.

\subsection{Type Erasure e Polimorfismo: i Metodi Ponte (Bridge Methods)}

La cancellazione del tipo si applica ovviamente anche ai metodi generici; generalmente, si tende a pensare al metodo:
\begin{center}
    \texttt{public static <T extends Comparable> T min(T[] a)}
\end{center}

come a una famiglia di metodi. In realtà, dopo la cancellazione del tipo, ciò che resta è il singolo metodo:
\begin{center}
    \texttt{public static Comparable min(Comparable[] a)}
\end{center}

La cancellazione del tipo sui metodi generici comporta però qualche complicazione quando una sottoclasse prova a fare l'override di un metodo ereditato da una classe generica. Consideriamo ad esempio: 

\begin{lstlisting}[language=Java]
    class DateInterval extends Pair<LocalDate>{
        public void setSecond(LocalDate second)
        {
            if (second.compareTo(getFirst()) >= 0)
                super.setSecond(second);
        }
        . . .
        }
\end{lstlisting}

Qui è stato eseguito l'override di \texttt{setSecond} per essere sicuri che la seconda data non sia mai precedente alla prima. Il problema è che nella classe \texttt{Pair}, dopo la cancellazione del tipo, esiste il metodo \texttt{setSecond(Object second)}, che viene ereditato da \texttt{DateInterval}; questo è un metodo diverso da \texttt{setSecond(LocalDate second)} che abbiamo definito (ha una firma diversa, come se si fosse fatto l'overload al posto di un override). 

Per evitare che il polimorfismo fallisca, il compilatore genera automaticamente un \textit{metodo ponte} nella sottoclasse:
\begin{lstlisting}[language=Java]
    public void setSecond(Object second){
         setSecond((LocalDate) second); 
         }
\end{lstlisting}

\begin{dettagli}[Dettagli: Logica di Rilevamento dei Metodi Ponte]

    Il compilatore genera un \textit{bridge method} quando rileva una discrepanza tra l'intento del programmatore e i vincoli della \textit{Type Erasure}.

    Il compilatore esegue i seguenti passaggi logici durante l'analisi della sottoclasse:

    \begin{enumerate}
        \item \textbf{Analisi delle Firme}: Identifica un metodo nella sottoclasse che corrisponde a un metodo della superclasse generica (es. \texttt{setSecond(LocalDate)} vs \texttt{setSecond(T)}).
        \item \textbf{Previsione della Cancellazione}: Calcola come appariranno le firme nel bytecode (es. \texttt{LocalDate} rimane tale, ma \texttt{T} diventa \texttt{Object}).
        \item \textbf{Rilevamento del Conflitto}: Se le firme cancellate non coincidono, il metodo non e' piu' un \textit{override} valido per la JVM.
        \item \textbf{Sintesi}: Genera automaticamente un metodo con la firma cancellata della superclasse che delega il lavoro al metodo della sottoclasse.
    \end{enumerate}

    \smallskip

    Il compilatore "capisce" la necessità del ponte perché possiede la tabella dei tipi completi prima che la cancellazione dei dati abbia luogo. Il ponte è il meccanismo che permette di mappare una chiamata generica \textit{raw} (basata su \texttt{Object}) alla specifica implementazione \textit{type-safe} della sottoclasse.
\end{dettagli}

La situazione diventa ancora più strana con i metodi che restituiscono valori; consideriamo ad esempio: 
\begin{lstlisting}[language=Java]
    class DateInterval extends Pair<LocalDate>{
        public LocalDate getSecond() { return (LocalDate) super.getSecond(); }
        . . .
        }
\end{lstlisting}

Nella classe \texttt{DateInterval} avremmo a questo punto due metodi con lo stesso nome \texttt{getSecond()} e nessun parametro: uno che restituisce un \texttt{Object} e uno che restituisce un \texttt{LocalDate}; come ben sappiamo, in Java questo sarebbe illegale, perché non si possono avere due metodi con la stessa firma che restituiscono tipi diversi. Tuttavia, nella Virtual Machine, un metodo è identificato univocamente dalla combinazione di nome, parametri \textbf{e tipo di ritorno}. Il compilatore può quindi generare bytecode che la JVM gestisce correttamente, anche se noi non potremmo scriverlo a mano. Il compilatore Java si comporta in questi casi quasi come un "hacker" che sfrutta le libertà della Virtual Machine per aggirare i limiti del linguaggio stesso.

\begin{warningbox}[Firma del Metodo: Java vs JVM]
    Esiste una distinzione fondamentale tra come il linguaggio Java e la JVM identificano i metodi:
    \begin{itemize}
        \item \textbf{Java Language}: Un metodo è identificato univocamente dalla sua \textit{signature} (nome e tipi dei parametri). Il tipo di ritorno viene ignorato per evitare ambiguità durante l'invocazione.
        \item \textbf{JVM Bytecode}: Un metodo è identificato dal nome, dai parametri e dal \textbf{tipo di ritorno}.
    \end{itemize}
    Questa differenza è ciò che permette al compilatore di sintetizzare i \textit{bridge methods}. Il compilatore genera due metodi con la stessa firma Java ma con ritorni diversi, permettendo alla JVM di risolvere correttamente le chiamate polimorfiche dopo la \textit{type erasure}.
\end{warningbox}

\begin{observation}[Ritorno Covariante e Bridge Methods]
    I \textit{bridge methods} non sono limitati ai tipi generici; vengono impiegati anche per gestire i \textbf{tipi di ritorno covarianti}. Come abbiamo visto nel capitolo su ereditarietà e polimorfismo, si ha covarianza quando una sottoclasse effettua l'override di un metodo specificando un tipo di ritorno più restrittivo rispetto a quello della superclasse.

    Consideriamo ad esempio:
    \begin{lstlisting}[language=Java]
        class Animal {
            public Animal getIdentity() { return this; }
        }

        class Cat extends Animal {
            @Override
            public Cat getIdentity() { return this; } // Ritorno covariante
        }
    \end{lstlisting}

    Nella classe \texttt{Cat}, il compilatore sintetizza effettivamente due metodi nel bytecode per garantire la compatibilità:
    \begin{itemize}
        \item \texttt{Cat getIdentity()}: Il metodo definito dal programmatore.
        \item \texttt{Animal getIdentity()}: Il \textit{bridge method} che sovrascrive il metodo di \texttt{Animal} e chiama internamente la versione specifica.
    \end{itemize}
\end{observation}







